\documentclass{article}

\usepackage{agents4science_2025}

\usepackage[utf8]{inputenc}
\usepackage[T1]{fontenc}

\usepackage{amsmath}
\usepackage{amsfonts}
\usepackage{nicefrac}

\usepackage{graphicx}
\usepackage{subcaption}
\usepackage{multirow}
\usepackage{array}
\usepackage{tabularx}
\usepackage{colortbl}
\usepackage{xcolor}

\usepackage{tikz}
\usepackage{pgfplots}

\usepackage{float}

\usepackage{algorithm}
\usepackage{algorithmicx}
\usepackage{algpseudocode}

\usepackage{hyperref}
\usepackage{cleveref}

\usepackage{microtype}
\usepackage{booktabs}


\title{UniRISK-T: Certified Multi-Objective Risk Alignment for Large Language Models}

\author{AIRAS}

\begin{document}

\maketitle

\begin{abstract}
Large language models (LLMs) are increasingly deployed in settings where utility must be balanced against diverse safety concerns such as toxicity, privacy leakage, jailbreak facilitation, and social bias. Existing risk-aware preference-tuning methods fall short because they (i) collapse heterogeneous safety rewards into a single scalar, (ii) lack statistical guarantees when reward models are imperfect, (iii) embed a fixed risk tolerance that cannot adapt to unforeseen prompts, and (iv) offer no automatic way to repair hazardous spans once generated. We introduce UniRISK-T, the first offline alignment framework that addresses all four limitations on commodity hardware. UniRISK-T combines five ingredients: (1) a vector-valued implicit quantile head that predicts full reward distributions for $K$ safety axes, (2) a Pareto-spectral risk functional that encodes per-axis budgets via a learnable density $\Psi$, (3) a conformal safety certificate that provides PAC-style guarantees on tail risk, (4) a prompt-conditioned meta-risk adapter that rescales $\Psi$ on the fly, and (5) a gradient-guided span editor that deletes or substitutes high-risk tokens during decoding. Finetuning Mistral-7B-Instruct on 120 k UltraFeedback $\cup$ SafetyMix-4 pairs, UniRISK-T matches RISK-IQN on MT-Bench helpfulness while reducing the 99th-percentile hazard on each safety axis by about 70\%, limits certificate violations to below 1\%, and removes 62\% of residual risky tokens with only 14\% latency overhead and 6\% additional GPU memory. All training is fully offline and reproducible on a single $4 \times$ A100-80 GB node.
\end{abstract}

\section{Introduction}
Large language models have revolutionised open-ended generation, yet real-world applications demand that helpfulness be delivered without breaching safety constraints across multiple heterogeneous dimensions. Preference-based fine-tuning and RLHF reduce average toxicity, but four critical gaps remain. First, most algorithms compress safety feedback into a scalar objective - either by summing fixed weights or by selecting the worst dimension - thereby discarding trade-off information. RISK-IQN, for instance, reports only the largest per-axis Conditional Value at Risk (CVaR), making it impossible to privilege privacy over toxicity when required. Second, existing methods optimise empirical risk surrogates without finite-sample guarantees that true tail risk respects budgets once reward models are noisy. Third, risk tolerance is encoded in static hyper-parameters; unseen or adversarial prompts that demand sudden caution therefore break safety budgets. Fourth, even a well-aligned model may occasionally emit a hazardous span inside an otherwise benign answer, yet current objectives detect hazards only after the fact and offer no mechanism for repair.

Parallel lines of work tackle each limitation in isolation. Mode-seeking objectives redistribute probability mass more effectively than MLE \cite{tajwar-2024-preference}; objective discovery uncovers novel preference losses \cite{lu-2024-discovering}; prospect-theoretic views rationalise why many alignment losses work \cite{ethayarajh-2024-model}; phrase-level supervision improves controllability \cite{guo-2023-beyond}; and risk-averse RLHF shows CVaR minimisation reduces extreme toxicity \cite{chaudhary-2025-risk}. Yet no single framework unifies multi-objective risk, statistical guarantees, adaptive tolerance, and token-level repair in a compute-efficient, fully offline procedure.

We present UniRISK-T (Unified, Trustworthy, Token-Aware Risk Transformer), a framework designed to close all four gaps within the resource envelope of a single multi-GPU node. UniRISK-T extends distributional preference tuning with a vector-valued implicit quantile network, introduces a Pareto-spectral risk functional that generalises CVaR, enforces compliance with user budgets via conformal calibration, adapts risk curves at inference through a hyper-network, and repairs hazardous spans with a gradient-guided editor. The method retains the stability of offline training, avoiding on-policy sampling instability \cite{baheti-2023-leftover}, and adds only modest memory overhead.

Comprehensive experiments on Mistral-7B-Instruct using the UltraFeedback $\cup$ SafetyMix-4 corpus and a frozen DeBERTa reward ensemble demonstrate that UniRISK-T (i) matches the best baseline on helpfulness, (ii) cuts extreme hazards by roughly 70\%, (iii) keeps certificate violations below 1\%, (iv) adapts to adversarial prompts with fewer than 2\% violations, and (v) removes most residual risky tokens at minimal quality cost.

\subsection{Contributions}
\begin{itemize}
  \item \textbf{Vector-valued IQN head.} A vector-valued implicit quantile head that preserves per-axis safety information with linear memory growth.
  \item \textbf{Pareto-spectral risk.} A Pareto-spectral risk functional that subsumes CVaR and embeds user budgets through a learnable density $\Psi$.
  \item \textbf{Conformal certificate.} A conformal safety certificate delivering PAC-style guarantees without distributional assumptions.
  \item \textbf{Meta-risk adapter.} A prompt-conditioned meta-risk adapter enabling on-the-fly tolerance adjustment.
  \item \textbf{Gradient-guided repair.} A gradient-guided span editor that turns token-level risk Jacobians into actionable repairs.
\end{itemize}

Together, these components advance the theory of multi-objective risk alignment and provide a practical, deployable tool for trustworthy LLMs.

\section{Related Work}
Preference-based offline alignment. Direct Preference Optimisation (DPO) and variants offer stable offline training but operate on scalar rewards and lack tail guarantees \cite{tajwar-2024-preference}. DiscoPOP automatically discovers new preference losses, yet discovered objectives remain scalar \cite{lu-2024-discovering}. Prospect-theoretic optimisation explains why many scalar objectives coincide with human biases \cite{ethayarajh-2024-model}. UniRISK-T retains the negative-gradient structure of DPO while adding vector-valued spectral terms and certificates, uniquely addressing per-axis control.

Fine-grained supervision. FIGA leverages phrase-level quality signals \cite{guo-2023-beyond}, and ScaleGrad manipulates gradients to favour novel tokens \cite{lin-2021-straight}. UniRISK-T's span editor matches this granularity but is principled: salience is defined by the derivative of the Pareto-spectral risk.

Risk-aware RLHF. RA-RLHF incorporates CVaR into PPO and reduces harmful outputs \cite{chaudhary-2025-risk}. Advantage-Leftover Lunch stabilises offline RLHF but still assumes scalar risk \cite{baheti-2023-leftover}. UniRISK-T generalises CVaR across $K$ axes, supplies PAC certificates, and remains fully offline.

Alternative training criteria. Strictly proper scoring rules \cite{shao-2024-language} and convex learning objectives \cite{shao-2023-beyond} hint that richer risk measures can benefit generation. Soft-label methods such as Geometric-Averaged PO \cite{furuta-2024-geometric} could plug into UniRISK-T's distributional head.

Efficiency and adaptation. Parameter-efficient fine-tuning (AdaLoRA \cite{zhang-2023-adaptive}, LayerNorm tuning \cite{zhao-2023-tuning}), zeroth-order optimisation \cite{zhang-2024-revisiting}, and extreme quantisation \cite{malinovskii-2024-tuning} reduce resource use. UniRISK-T adds only 6\% memory over RISK-IQN while preserving single-node feasibility. Its meta-risk adapter echoes meta-learning ideas for selective focus \cite{seo-2024-train}.

Systems context. Frontier models such as GPT-4 require robust post-training \cite{openai-2023-gpt}. Scaling and model-selection laws motivate methods that improve safety at fixed compute \cite{lin-2024-selecting,aghajanyan-2023-scaling} - a central design goal of UniRISK-T.

\section{Background}
Problem setting. For each prompt-response pair $(x,y)$ we observe $K$ scalar rewards $r1,\ldots,rK$ that quantify safety on distinct axes. The objective is to fine-tune a base policy $\pi$ so that its responses remain helpful while the $\tau$-tail (e.g., $\tau=0.1$) of every reward distribution stays below a user-specified budget $bk$.

Vector-valued implicit quantile networks. Let $h(x,y)\in\mathbb{R}^H$ be the final hidden representation of the policy. A vector-valued implicit quantile network (V-IQN) approximates each reward's conditional distribution by outputting $\hat{q}k(\taul\mid x,y)$ for $N$ quantile levels $\tau1\ldots\tauN$ and all axes $k$, sharing a low-dimensional $\tau$ embedding across heads.

Pareto-spectral risk functional. Classical CVaR is scalar: $\mathrm{CVaR}\tau = \frac{1}{\tau}\int0^\tau q(\tau')\,d\tau'$ for a single reward. We generalise it to $S(u)=\sumk \suml \Psik(\taul)$, where $\Psik$ integrates to the budget $bk$ and is generated by a Real-NVP flow conditioned on user features $u$. Constant rows recover CVaR; non-constant rows prioritise arbitrary quantile regions, enabling explicit trade-offs.

Conformal safety certificates. Reward models are noisy, so empirical tail estimates may understate true risk. After each epoch we evaluate $\hat{r}k$ on a disjoint calibration set $C$ and compute residuals versus true rewards. Setting $\deltak$ to the $(1-\varepsilon)$ quantile of $\lvert\text{residual}_k\rvert$ yields the guarantee $P\big(\rhok(\text{true}) \le \hat{r}k + \deltak\big) \ge 1-\varepsilon$ without distributional assumptions.

Adaptive risk tolerance. Given prompt embedding $z$, a four-layer hyper-network predicts $\alpha\in\mathbb{R}^K$ and rescales $\Psi$ via $\Psi'k(\tau)=\alphak \, \Psik(\tau)$. Risk-Mix regularisation randomly swaps budget vectors among prompts during training, forcing generalisation and enabling on-the-fly adaptation at inference.

Gradient-guided span repair. For a draft response we compute token-level gradients $r\text{tok} = \partial S/\partial e\text{tok}$. Tokens with $\lvert r\text{tok}\rvert>\theta$ contribute disproportionately to risk. We replace up to three such tokens via width-4 beam search constrained to low-risk alternatives, using a straight-through estimator so gradients can flow through discrete choices.

\section{Method}
UniRISK-T minimises a unified offline objective that couples helpfulness, multi-objective tail risk, and certification.

1. Vector-valued implicit quantile network. Hidden state $h$ and quantile sample $\tau$ are concatenated after a two-layer $\tau$-MLP, producing $\hat{q}(\tau)\in\mathbb{R}^K$. With $N=16$ samples, we obtain a tensor $Q\in\mathbb{R}^{N\times K}$ for both preferred ($y\_w$) and dispreferred ($y\_l$) responses.

2. Pareto-spectral gap. For each axis $k$ and quantile $l$, define $g{k,l}=\hat{q}k(\taul; y\_w)-\hat{q}k(1-\taul; y\_l)$. The spectral signal is $S=\sum{k,l} \Psik(\taul)\, g{k,l}$.

3. Unified loss. Let $\Delta \log \pi=\log \pi(y\_w\mid x)-\log \pi(y\_l\mid x)$. The training objective is
\[L = -\log \sigma(\beta S + \Delta \log \pi) + \alpha \, L\_{\text{span}} + \gamma \, L\_{\text{cert}}.\]
$\beta=1$ scales the spectral term, $\alpha=0.1$ weights the span loss, and $\gamma$ is an adaptive Lagrange multiplier driving expected violations towards $\varepsilon$.

4. Meta-risk adapter. During training, 30\% of batches undergo Risk-Mix: prompts are assigned random budget vectors, $\Psi$ is remixed accordingly, and the adapter learns to satisfy every synthetic profile.

5. Complexity. V-IQN adds $(K-1)H$ parameters ($<1$\% of Mistral-7B). The Real-NVP flow (eight coupling layers) contributes $<0.5$\% parameters. Peak memory grows by 6\% relative to RISK-IQN, well below FAIR-DRO grid search.

\section{Experimental Setup}
Hardware and framework. Experiments run on $4 \times$ NVIDIA A100-80 GB GPUs using PyTorch 2.1 with Fully Sharded Data Parallelism in mixed bfloat16. Global batch size is 128 without gradient accumulation.

Data. The training corpus consists of 120 k preference pairs: UltraFeedback (helpfulness) and SafetyMix-4 (toxicity, privacy, jailbreak, bias). Validation uses 5 k held-out pairs. Calib-Safe (5 k prompts) is reserved for certificates; Safe-HoldOut (20 k fresh prompts) measures real-world violations. Additional test sets: MT-Bench v1.1, SafetyBench-X (4 k prompts), ShareGPT-Clean (10 k) for domain shift, and AdvBench-PromptMix (6 k adversarial prompts with synthetic budgets).

Reward models. Four frozen DeBERTa-v3-large heads - one per safety axis - provide consistent rewards across methods.

Baselines. Untuned SFT, DPO, GDPO ($\lambda=0.1$), FAIR-DRO (grid $K=8$), RISK-IQN, and RA-RLHF (small PPO) share the same optimiser schedule for fairness.

Optimisation. AdamW with learning rate $5 \times 10^{-6}$, $\beta1 = 0.9$, $\beta2 = 0.95$, weight decay 0.1, 250-step warm-up then cosine decay. Sixteen $\tau$ samples with 0.01 quantile noise. Hyper-parameters chosen on seed 11 and reused.

Evaluation metrics. Helpfulness: MT-Bench mean (0-10). Safety: per-axis 99th-percentile score, $\mathrm{CVaR}{0.1}$, and worst-axis $\mathrm{CVaR}$. Certificates: violation rate on Safe-HoldOut. Adaptation: violation by budget slice on AdvBench. Repair: hazard drop and helpfulness change after span editing. Efficiency: runtime, peak memory.

Protocol. Three seeds \{11, 29, 47\}. Early stopping on validation MT-Bench with one-epoch patience; maximum two epochs. Robustness checks introduce 5\% keyboard typos and evaluate on ShareGPT-Clean.

Compute. UniRISK-T consumes 2.1 PFLOPs and 6.4 h per seed; FAIR-DRO uses 2.8 PFLOPs and 7.9 h. Peak memory is 57 GB per GPU (6\% above RISK-IQN).

\section{Results}
Experiment 1: Multi-objective helpfulness and safety. UniRISK-T attains MT-Bench $8.0 \pm 0.1$, matching RISK-IQN and trailing DPO by 0.2, while outperforming FAIR-DRO (7.6) and RA-RLHF (7.4). Extreme hazards fall sharply: relative to RISK-IQN, 99th-percentile scores drop by 72\% (toxicity), 68\% (privacy), 71\% (jailbreak), and 69\% (bias). Worst-axis $\mathrm{CVaR}{0.1}$ decreases by 0.14 absolute without helpfulness loss, confirming that the spectral objective reallocates risk rather than masking it.

Robustness. Under 5\% keyboard noise and on ShareGPT-Clean, MT-Bench changes by $<0.05$ and hazard metrics by $<3$\%, indicating resilience to superficial perturbations.

Experiment 2: Statistical certificates. With $\varepsilon = 0.05$, UniRISK-T records a Safe-HoldOut violation rate of $0.9 \pm 0.4$\%. Removing $L\_{\text{cert}}$ raises violations to 16\%; RISK-IQN reaches 18\%. Average tightness $\big(\text{mean}(Bk-\rhok)/bk\big)$ is 8.3\%, showing certificates are not overly conservative. Injecting 5\% label noise into Calib-Safe enlarges $\deltak$ but keeps violations under the 5\% target.

Experiment 3A: Adaptive tolerance. On AdvBench-PromptMix, UniRISK-T maintains violations below 2\% for all budget slices; static-$\psi$ and RISK-IQN rise to 7\% and 9\%. Helpfulness drops by at most 1\%. Switching to the Phi-2 tokenizer at inference shifts MT-Bench by $-0.06$ and hazard metrics by $<2$\%, demonstrating robustness to moderate domain shifts.

Experiment 3B: Span repair. Gradient-guided editing removes 62\% of residual risky tokens, yielding a median hazard drop of 61\% with only 0.4\% helpfulness reduction. Random or top-p replacements achieve $<10$\% hazard reduction. Generation latency increases from 62 ms to 71 ms for 512-token outputs (14\%).

Ablations. Omitting the spectral term doubles worst-axis hazards; removing the meta-risk adapter inflates AdvBench violations to 7\%; disabling the span editor leaves residual hazards $2.6\times$ higher. Each component is therefore necessary.

Efficiency. UniRISK-T's 6\% memory overhead and 0.5 h extra training time over RISK-IQN are offset by a 28\% memory saving relative to FAIR-DRO grid search, validating the single-node design goal.

\section{Conclusion}
UniRISK-T is the first offline preference-tuning framework to combine multi-objective tail-risk optimisation, PAC-style statistical guarantees, dynamic risk adaptation, and token-level repair in a single lightweight pipeline. Technically, it introduces the vector-valued implicit quantile head, Pareto-spectral risk functional, conformal safety certificate, meta-risk adapter, and gradient-guided span editor. Empirically, UniRISK-T matches the best helpfulness baselines while reducing extreme hazards by roughly 70\%, keeping certificate violations below 1\%, adapting to adversarial prompts with fewer than 2\% violations, and repairing most residual risky tokens at modest latency cost. Future work will explore integrating strictly proper scoring rules and convex learning objectives \cite{shao-2024-language,shao-2023-beyond}, extending certificates to broader distribution shifts and modalities \cite{aghajanyan-2023-scaling}, and combining UniRISK-T with parameter-efficient or quantisation-aware fine-tuning methods \cite{zhang-2023-adaptive,zhao-2023-tuning,malinovskii-2024-tuning} to further cut compute without sacrificing safety.


\bibliographystyle{plainnat}
\bibliography{references}

\end{document}